% from ega4-11-fr.tex
\oldpage[IV-3]{163}
$\rho_{\mathcal G'_\lambda}:\mathcal G'_\lambda\to
(j'_{\lambda\mu})_*((j'_{\lambda\mu})^*(\mathcal G'_\lambda))$
correspondant à l'injection canonique
$j'_{\lambda\mu}:Y'_{\lambda\mu}\to Z'_\lambda$, comme on l'a vu dans
(11.9.5). L'assertion résulte alors des raisonnements de (11.9.2) et (11.9.3),
$Y_{\lambda\mu}$ et $v_{\lambda\mu}$ étant remplacés par
$Y'_{\lambda\mu}$ et $v'_{\lambda\mu}$.

On a une réduction semblable lorsque l'on veut prouver que la famille
$(u'_\lambda)$ est \emph{séparante} ($X$ et $X'$ étant affines)~: cela résulte
encore de (11.9.2) et (11.9.3).
\end{env}

\begin{proposition}[11.9.7]
\label{IV.11.9.7-fr}
Avec les notations de (11.9.1) et (11.9.5), supposons
$X=\operatorname{Spec}(A)$ et $X'=\operatorname{Spec}(A')$ affines, et
supposons en outre que $A'$ soit un $A$-module projectif. Alors, si la famille
$(u_\lambda)$ est séparante, toute section $t'$ de $\mathcal F'$ au-dessus de
$X'$ dont les images par tous les $u'_\lambda$ sont nulles, est elle-même
nulle.
\end{proposition}

\begin{proof}
On a vu (11.9.6) qu'on peut se borner au cas où tous les $Z_\lambda$ sont égaux
à $X$. La proposition est alors conséquence du lemme suivant~:
\end{proof}

\begin{lemma}[11.9.7.1]
\label{IV.11.9.7.1-fr}
Soient $A$ un anneau, $(M_\lambda)_{\lambda\in L}$ une famille de $A$-modules,
$M$ un $A$-module et pour chaque $\lambda$,
$u_\lambda:M\to M_\lambda$ un homomorphisme. Supposons que l'intersection des
noyaux des $u_\lambda$ soit réduite à 0. Alors, pour tout $A$-module projectif
$N$, l'intersection des noyaux des homomorphismes
$u_\lambda\otimes 1:M\otimes_A N\to M_\lambda\otimes_A N$ est réduite à 0.
\end{lemma}

\begin{proof}
En effet, $N$ est facteur direct d'un $A$-module libre $L$, et il suffit
évidemment de prouver que l'intersection des noyaux des homomorphismes
$v_\lambda:M\otimes_A L\to M_\lambda\otimes_A L$ est réduite à 0, puisque
$u_\lambda\otimes 1:M\otimes_A N\to M_\lambda\otimes_A N$ est la restriction
de $v_\lambda$. Mais l'assertion résulte alors trivialement de l'hypothèse.
\end{proof}

\begin{remark}[11.9.8]
\label{IV.11.9.8-fr}
Nous ignorons si, sous les hypothèses de la proposition (11.9.7), la famille
$(u'_\lambda)$ est \emph{séparante}~: il faudrait en effet (11.9.4) prouver
qu'une section $t'$ de $\mathcal F'$ \emph{au-dessus d'un ouvert}
$D(h')\subset X'$ (où $h'\in A'$) telle que les $u'_\lambda(t')$ soient toutes
nulles, est elle-même nulle. Or, on ne peut appliquer la proposition (11.9.4)
à $D(h')=\operatorname{Spec}(A'_{h'})$, car du fait que $A'$ soit un
$A$-module projectif (même libre), il ne résulte pas que $A'_{h'}$ soit un
$A$-module projectif. Par exemple, on peut prendre pour $A$ un anneau de
valuation discrète, pour $A'$ un anneau de valuation discrète qui soit un
$A$-module libre de rang 2, et pour $A'_{h'}$ le corps des fractions de $A'$.
\end{remark}

On a toutefois le résultat suivant~:

\begin{corollary}[11.9.9]
\label{IV.11.9.9-fr}
Soient $X$ un préschéma artinien, $g:X'\to X$ un morphisme plat (on notera que
ces deux conditions sont satisfaites si $X$ est le spectre d'un \emph{corps}
et $g$ un morphisme \emph{quelconque}). Alors, avec les notations de (11.9.1)
et (11.9.5), si la famille $(u_\lambda)$ est séparante, il en est de même de la
famille $(u'_\lambda)$.
\end{corollary}

\begin{proof}
On peut évidemment se borner au cas où $X=\operatorname{Spec}(A)$ est le
spectre d'un anneau local artinien $A$ (\textbf I, 6.2.2)~; on note ensuite que
pour tout ouvert affine $U'=\operatorname{Spec}(A')$ de $X'$, $A'$ est un
$A$-module plat, donc \emph{projectif} $(\mathbf 0_{\mathrm{III}},10.1.3)$.
Il suffit donc d'appliquer (11.9.7) à tout ouvert affine de $X'$ pour obtenir
le corollaire.
\end{proof}

\begin{theorem}[11.9.10]
\label{IV.11.9.10-fr}
Soient $X$ un préschéma, $(u_\lambda)$ une famille d'homomorphismes
(11.9.1.1), $g:X'\to X$ un morphisme, $(u'_\lambda)$ la famille
d'homomorphismes déduite de $(u_\lambda)$ par le changement de base $g$
(11.9.5).
\begin{enumerate}[label=(\roman*)]
  \item Si $g$ est un morphisme fidèlement plat et si la famille
  $(u'_\lambda)$ est séparante, alors la famille $(u_\lambda)$ est séparante.
  